For this question use sky map Map-OP01. This is an angle preserving
projection of the sky.

The sky projected on the planetarium dome is as seen from the solar system
barycentre. The Earth's orbit as seen from this position has been marked in
green both on the projected sky and in sky map Map-OP01. The first 3 minutes
are for dark adaptation and familiarization of the sky.

After first 3 minutes, four Trans-Neptunian Objects (TNOs) labelled O1, O2,
O3 and O4 will begin to appear on the dome. The brightness of the TNOs has
been enhanced to make them visible. No other solar system object is visible
in the projected sky. Time has been sped up to make their movements
noticeable.

We will assume that the orbital plane of a hypothesized new planet has an
inclination equal to the mean of the inclinations of the orbits of these 4
TNOs.

% FIGURE NEEDED: sky map Map-OP01 (angle preserving projection with the
% Earth's orbit marked in green) — distributed as a separate chart with the
% exam; not part of the question paper PDF.

\begin{description}
\item[(OP01.1)] Arrange the 4 TNOs in ascending order of the lengths of their
  semi-major axes and write their names in the appropriate boxes in the
  Summary Answersheet. Show measurements in the working sheet to justify your
  answer. \hfill[8]
\item[(OP01.2)] Trace at least 50\% of the visible trajectories for each of
  the 4 TNOs on the map Map-OP01. From these traces find the inclination
  angles, $i_1, i_2, i_3, i_4$, of the 4 TNOs, respectively, with respect to
  the Ecliptic. Hence obtain the best estimate of the inclination,
  $i_{\mathrm{hp}}$, of the plane of the orbit of the hypothesized planet.
  \hfill[17]
\end{description}
